\section{Dynamics, Flows, and External Comparison}\label{sec:dynamics}

\subsection{Dynamics}

The binary post coefficient is not an event-study design.  I estimate quarterly coefficients for all four nonreference quintiles, omitting 2022Q4 (October--November; December is the transition month).  The Q5 path has 23 preperiod and 15 postperiod coefficients in pooled and family-month specifications.

Figure~\ref{fig:dynamics} shows the rebuilt Q5 path.  Relative to the omitted
2022Q4 coefficient, joint tests reject that all 23 preperiod coefficients are
zero in the pooled ($p=.0171$) and family-conditioned ($p=.0053$) models.  A
reference-invariant test of equality within the entire preperiod rejects for
the pooled path ($p=.0121$) but not for the family-conditioned path
($p=.0641$).  The appendix reports fixed pre-pandemic, post-reopening, and
pandemic-exclusion windows.  An insignificant anchored linear slope cannot
override the unrestricted pooled test, while the family result illustrates
the importance of the unusually short 2022Q4 reference.  Neither pattern
supports a causal January-2023 interpretation.
When 2020Q2--2020Q4 are excluded, the pooled test no longer rejects at five
percent under either the short-reference null ($p=.0772$) or the
reference-invariant within-window null ($p=.0676$).  For the family-conditioned
path the corresponding values are $p=.0079$ and $p=.1048$, underscoring that
the conclusion depends on the null as well as the calendar window.

\begin{figure}[!htbp]
\centering
\includegraphics[width=0.94\textwidth]{r3_figure2_dynamics.pdf}
\caption{Quarterly Q5--Q1 young-relative employment coefficients}\label{fig:dynamics}
\tabnote{Coefficients are relative to Q1 and omitted 2022Q4 under rebuilt labels; lines are pointwise 95-percent occupation wild-score intervals.  The online appendix reports simultaneous bands and the Q2--Q5 system.  The conditioned model gives each SOC2 family its own young-relative monthly path.}
\end{figure}

Three related targets separate a genuine estimand comparison from an event-time
normalization artifact.  The static grouped-binomial coefficient $S$ is
$-0.1321$ pooled and $-0.0217$ after family conditioning.  The calendar-
weighted post functional relative to omitted 2022Q4, $P$, is $-0.1199$
($[-0.2637,0.0239]$) and $-0.2074$ ($[-0.4252,0.0103]$).  By contrast, the
reference-invariant post-minus-pre functional $D$ is $-0.1314$ and $-0.0217$.
The paired $D-S$ differences are 0.0007 ($p=.454$) and less than 0.0001 in
absolute value ($p=.995$).  Thus $D$ nearly reproduces the static estimates;
the design does not detect the family-model $P-S$ difference at five percent
($-0.1858$, $p=.074$).  What it does detect is that family conditioning changes
$P-S$ by $-0.1980$ ($p=.021$), while it does not detect a change in $D-S$
($-0.0008$, $p=.693$).  The short-reference functional therefore responds to
conditioning differently from the reference-invariant comparison.  This
numerical reconciliation is not an equivalence claim and does not make any
estimand causal.

Official \texttt{HonestDiD} 0.2.8 applies smoothness and relative-magnitude restrictions to the full event vector and covariance for $P$, not for $S$ or $D$.  Both the full 2017Q1--2022Q3 calibration and a consecutive 2021Q1--2022Q3 post-reopening calibration are reported.  Under the latter, the family-conditioned smoothness interval is $[-0.772,-0.007]$ when differential-trend bias is restricted to be exactly linear ($M=0$), but expands to $[-1.100,0.207]$ at the first positive grid value, $M=0.005$.  Both conventional $P$ intervals already include zero, so no positive zero-exclusion threshold exists.  The exercise cannot repair rejected preperiod equality or identify adoption.

Onset estimates from November 2022 through June 2023 are similar within each structure, but common timing does not distinguish AI from coincident shocks.  Ending in December 2024 gives a pooled coefficient of $-0.111$ ($[-0.201,-0.021]$); its paired movement relative to the July-2026 estimate is $0.021$ ($[-0.007,0.049]$).  The corresponding movements at the September- and December-2025 endpoints are 0.0225 and 0.0212, essentially the same magnitude despite different interval endpoints.  The full window remains the canonical target and December 2024 a prominent pre-production-break benchmark; neither is causally clean.  Splitting the postperiod instead gives $-0.111$ in 2023--24 and $-0.168$ in 2025--26; their paired difference is $-0.057$ ($[-0.125,0.012]$), so the official-weight design does not detect a change between eras.  Family-conditioned endpoint changes remain near zero, and survey-gap exclusions are small.  A coding-stable post-2020 fit gives $-0.118$ pooled and $-0.030$ with family-month effects; their paired movements from the full-window models are $0.014$ ($[-0.029,0.057]$) and $-0.009$ ($[-0.105,0.088]$).  Quintile- and occupation-specific month-of-year controls move either coefficient by at most $0.0011$.  All six replacement fits pass the same numerical certification; the earlier failed attempts remain in the failure registry rather than being treated as results.

\subsection{Linked CPS margins}

Table~\ref{tab:flows} estimates distinct adjacent- and twelve-month margins using validated links and official longitudinal weights.  Every primary interval includes zero.  Normal-theory MDE$_{80}$ values range from 0.115 to 0.336 log point, so these are low-resolution nondetections rather than evidence of absent flow responses.

\begin{table}[!htbp]
\centering
\caption{Corrected CPS flow margins}\label{tab:flows}
\resizebox{\textwidth}{!}{%
\begin{tabular}{lrrrr}
\toprule
Endpoint and margin & Raw link records at risk / entries & Coefficient & 95\% interval & Normal MDE$_{80}$ \\
\midrule
Adjacent employment exit & 3,346,227 & $0.132$ & $[-0.042,0.307]$ & 0.254 \\
Adjacent reported-occupation change & 3,207,598 & $0.003$ & $[-0.118,0.123]$ & 0.174 \\
Adjacent entry-destination allocation & 96,981 & $-0.070$ & $[-0.250,0.110]$ & 0.260 \\
Twelve-month employment exit & 1,399,376 & $0.123$ & $[-0.111,0.357]$ & 0.336 \\
Twelve-month reported-occupation endpoint change & 1,110,024 & $-0.018$ & $[-0.097,0.061]$ & 0.115 \\
Twelve-month entry-destination allocation & 82,573 & $-0.052$ & $[-0.251,0.147]$ & 0.282 \\
\bottomrule
\end{tabular}}
\tabnote{Counts are repeated endpoint-link records, not unique persons; an individual may contribute more than one eligible link.  Coefficients are post-January-2023 changes in the ages-22--25 versus ages-26--65, beta-Q5 versus beta-Q1 log rate ratio.  Entry destination is a log allocation ratio conditional on an observed entry, not an employment-finding or employer-hiring rate.  Primary intervals use 9,999 occupation wild-score draws and official longitudinal weights.  MDE$_{80}=2.80\widehat{SE}$ uses the analytic occupation-cluster standard error and is a normal-theory precision diagnostic, not the wild-score rejection threshold and not evidence that a smaller effect is absent.}
\end{table}


Employment exit starts from employed origins; reported occupation change requires employment at both endpoints; entry destination conditions on nonemployment-to-employment links and measures allocation across exposure groups.  The last is neither a finding probability nor employer hiring.  Twelve-month links are endpoints, not sums of monthly transitions.  Among fixed-support employed origins, postperiod twelve-month retention ranges from 48.7 to 52.8 percent across young quintiles and from 63.9 to 70.2 percent across older quintiles; preperiod ranges are 50.1--54.0 and 68.8--73.3 percent.  Hours ($-0.119$, $[-0.295,0.057]$) and conditional weekly earnings ($-0.0245$, $[-0.1024,0.0534]$) are also imprecise and condition on employment; the legacy earnings series ends in March 2023 and contains only three postperiod months, while a full-window
\texttt{EARNWEEK2} analysis remains unavailable in the authorized extracts.

Observable poststratification does little to the adjacent employment-exit
probability contrast: the descriptive Q5--Q1, young--older, post--pre difference
is 0.0088 under official link weights and 0.0086 after reweighting retained
links to the eligible-origin distribution.  At twelve months it moves from
0.0166 to 0.0210.  This reweighting requires outcome missingness to be
independent of linkage within the observed strata and includes a maximum
untrimmed factor of 102.7.  Without that assumption, the jointly feasible
linear accounting interval for twelve-month employment exit is
$[-1.570,1.575]$, leaving the sign of selection bias unrestricted.  These are
descriptive probability contrasts and missing-outcome bounds, not the nonlinear
regression coefficient and its confidence interval.

The annual primary rule excludes origins whose twelve-month link straddles the
January 2023 onset.  Coding instead the fraction of destination months exposed
gives an employment-exit coefficient of 0.158 ($[-0.070,0.387]$), compared
with the Table~\ref{tab:flows} primary estimate of 0.123
($[-0.111,0.357]$).  Both intervals contain zero;
no paired interval for their difference is available.

All flow exposure is defined at the origin where an origin occupation exists.  No current occupation is imputed to a nonemployed respondent.  Annual endpoints miss intervening jobs and transitions, and positive longitudinal-weight eligibility changes the population; neither horizon is a mechanical decomposition of the stock coefficient.  Link-rate and observable-balance differences diagnose selection but do not determine its direction, and alternative clustering addresses dependence rather than missing outcomes.

\subsection{Payroll bridge and exposure architecture}

BCC's August 2026 19-percent ``kept pace'' shortfall is a stock-growth comparison, not this regression coefficient or a hiring rate.  Complete public membership and tie rules are unavailable, so I use an equal-occupation, tie-preserving approximation and label it a bridge.  Its top-two-versus-bottom-three CPS coefficient is $-0.0720$ ($[-0.1223,-0.0218]$); family-month conditioning gives $-0.0168$ ($[-0.0808,0.0473]$).  Their paired movement is $0.0553$ ($[-0.0013,0.1118]$), so the design does not detect a difference.  Differences in universe, ages, outcome, endpoints, employer controls, and worker--firm flows remain.

That grouping bridge is distinct from changing the regression objective.  Giving each of the 468 occupations equal total weight, while retaining \texttt{WTFINL} for within-cell age shares, changes the pooled Q5-by-post coefficient from $-0.1321$ to $-0.0833$; the paired difference is $0.0488$ ($[-0.0604,0.1581]$).  The family-month coefficient changes from $-0.0217$ to $-0.0047$, with paired difference $0.0170$ ($[-0.1587,0.1927]$).  These intervals are nondetections, not evidence that the objectives are equivalent.

The bridge therefore asks whether a closer public grouping changes the CPS stock pattern; it does not calibrate CPS transitions to proprietary employer flows.  The required entry, exit, job-to-job, occupation-switching, aging, and denominator quantities are not jointly observed.

Finally, $E_1+\lambda E_2$ varies the weight on tasks requiring complementary software.  Here $E_1$ is the directly accelerable task share and $E_2$ is the share requiring complementary software.  Each row uses rebuilt weights and common draws, and $\lambda=.5$ exactly reproduces beta.  Native-tail coefficients range from $-0.103$ at $\lambda=0$ to $-0.156$ at $\lambda=1$; every comparison with beta includes zero, even though 61.6 percent of stock changes quintile at $\lambda=0$.  This is large classification movement with limited coefficient resolution, not evidence about deployed technology.

Five non-$\lambda$ definitions give pooled native-support tail coefficients from $-0.099$ (direct-ability AIOE) to $-0.014$ (the reversed OECD AI-skills gap).  Every alternative-minus-beta paired occupation interval contains zero, although the OECD lower endpoint is $-0.0004$.  On the common AIOE support, some family-conditioned implementation contrasts do exclude zero; the result is therefore sensitive to how one construct is mapped in that specification, without implying that the alternatives measure the same technology.

\begin{table}[!htbp]
\centering
\caption{Weighting direct and software-complement tasks}\label{tab:lambda}
\resizebox{\textwidth}{!}{%
\begin{tabular}{rrrr}
\toprule
$\lambda$ in $D+\lambda S$ & Q5--Q1 coefficient & Difference from $\lambda=.5$ [95\% CI] & Stock changing group (\%) \\
\midrule
0.00 & $-0.103$ & $0.029\;[-0.044,0.103]$ & 61.6 \\
0.25 & $-0.101$ & $0.031\;[-0.005,0.068]$ & 25.5 \\
0.50 & $-0.132$ & --- & 0.0 \\
0.75 & $-0.143$ & $-0.011\;[-0.039,0.017]$ & 10.2 \\
1.00 & $-0.156$ & $-0.024\;[-0.064,0.017]$ & 25.5 \\
\bottomrule
\end{tabular}}
\tabnote{All rows use 468 occupations, reconstructed preperiod weights, tie-preserving native groups, Webb software conditioning, and 9,999 common occupation draws.  ``Stock changing group'' is preperiod employment in occupations whose quintile differs from beta ($\lambda=.5$).  The coefficient path is descriptive of the capability boundary; no paired beta comparison detects a difference at five percent.}
\end{table}

